\section{答题时间整体规划与答题技巧}

\subsection{试卷结构速览}

\begin{center}
\begin{tabular}{c|c|c|c}
\textbf{题型} & \textbf{题量} & \textbf{分值} & \textbf{建议用时} \\ \hline
单选题 & 7 题 & 28 分 & 15 分钟 \\
多选题 & 3 题 & 15 分 & 10 分钟 \\
实验题 & 2 题 & 15 分 & 12 分钟 \\
计算题 & 3 题 & 42 分 & 33 分钟 \\
检查 & --- & --- & 5 分钟 \\ \hline
总计 & 15 题 & 100 分 & 75 分钟
\end{tabular}
\end{center}

\subsection{答题顺序策略}

\textbf{原则}：先易后难，先熟后生，遇难则跳。

\textbf{推荐答题顺序}：

\begin{enumerate}
  \item \textbf{单选题 1--5 题}（基础题，8--10 分钟，确保全对）
  \item \textbf{多选题 1--2 题}（中档题，5--6 分钟，宁少勿滥）
  \item \textbf{实验题（第一道）}（基础实验，5--6 分钟）
  \item \textbf{计算题 1}（热学/光学/基础力学，8 分钟，必拿满分）
  \item \textbf{计算题 2}（力电综合，12 分钟，力争全对）
  \item \textbf{单选题 6--7 题}（中档偏难，5--6 分钟）
  \item \textbf{多选题第 3 题}（难题，4 分钟，只选确定选项）
  \item \textbf{实验题（第二道）}（创新实验，6 分钟）
  \item \textbf{计算题 3}（压轴题，第一问必拿，第二问尽力，13 分钟）
  \item \textbf{剩余时间检查}
\end{enumerate}

\textbf{底线}：75 分钟内，基础题（约 60 分）应确保拿到 55 分以上，中档题（约 25 分）力争拿 15 分，难题（约 15 分）争取步骤分。

\subsection{单选题答题技巧}

\begin{enumerate}
  \item \textbf{直接法}：概念清晰、公式可直接套用时使用，注意单位换算和数值估算
  \item \textbf{排除法}：
    \begin{itemize}
      \item 单位排除：结果的量纲与选项不一致时直接排除
      \item 极端值排除：取边界值（$0,\infty$）代入检查合理性
      \item 对称性排除：利用物理情境的对称性判断
    \end{itemize}
  \item \textbf{量纲分析法}：
    \begin{itemize}
      \item 物理量的单位可以辅助判断公式正误
      \item 例：$F=ma$ 的单位是 $\text{kg}\cdot\text{m/s}^2$，如果选项出现 $\text{kg}\cdot\text{m/s}$ 则必错
    \end{itemize}
  \item \textbf{估算与数量级}：
    \begin{itemize}
      \item 天题目中数据往往设计为整数或简单分数，不必精确计算
      \item 注意数量级：$10^{-3}$、$10^6$ 等容易出错的地方
    \end{itemize}
  \item \textbf{图象分析法}：
    \begin{itemize}
      \item $v-t$ 图看加速度和位移，$x-t$ 图看速度
      \item $U-I$ 图看电源电动势和内阻
      \item 图象斜率、截距、面积都有物理意义
    \end{itemize}
  \item \textbf{特值法}：
    \begin{itemize}
      \item 取特殊角度（$0^\circ,30^\circ,45^\circ,60^\circ,90^\circ$）
      \item 取特殊位置（最高点、最低点、平衡位置）
      \item 取特殊情况（光滑 $\to$ 无摩擦，轻杆 $\to$ 质量不计）
    \end{itemize}
\end{enumerate}

\subsection{多选题答题技巧}

\textbf{计分规则}：全部选对得 5 分，部分选对得 3 分，有错选得 0 分。

\textbf{核心策略——宁少勿滥}：
\begin{itemize}
  \item 只选 \textbf{100\% 确定} 的选项
  \item 不确定的选项 \textbf{宁愿不选}（多得 3 分 vs 丢 5 分）
  \item 尤其注意"可能"与"一定"的区分
\end{itemize}

\textbf{常见物理多选题特征}：
\begin{itemize}
  \item 图象类：$v-t$、$a-t$、$F-t$ 图结合运动分析
  \item 功能关系类：判断动能、势能、机械能如何变化
  \item 电磁感应类：判断电流方向、安培力方向、能量转化
  \item 电路动态分析：判断各电表示数变化、功率变化
  \item 天体运动类：比较线速度、角速度、周期、加速度
\end{itemize}

\subsection{实验题答题技巧}

\textbf{力学实验常见考点}：
\begin{itemize}
  \item 打点计时器纸带处理：$v_n = \dfrac{x_n+x_{n+1}}{2T}$，$a = \dfrac{\Delta x}{T^2}$
  \item 验证牛顿第二定律：平衡摩擦力方法、$a-F$ 图象不过原点原因
  \item 验证机械能守恒：$mgh = \frac12 mv^2$，误差来源（阻力）
  \item 探究平抛运动：频闪照片分析，$v_0 = \dfrac{x}{t}$
  \item 验证动量守恒：碰撞前后速度的测量方法
\end{itemize}

\textbf{电学实验常见考点}：
\begin{itemize}
  \item 伏安法测电阻：内外接法选择、滑动变阻器分压/限流
  \item 测电源电动势和内阻：$U-I$ 图象法，纵截距为 $E$，斜率为 $r$
  \item 多用电表使用：欧姆调零、读数、倍率选择
  \item 电表改装：串联分压（电压表）、并联分流（电流表）
  \item 传感器：光敏电阻、热敏电阻的特性
  \item 螺旋测微器和游标卡尺读数（常考第一空）
\end{itemize}

\textbf{实验题答题规范}：
\begin{itemize}
  \item 连线题：先串后并，电流从正极出发
  \item 读数题：注意估读位数（螺旋测微器读到 $0.001\,\text{mm}$）
  \item 作图题：描点后用直尺画直线，注意标度
  \item 表达式推导：先写原始公式再代入推导
\end{itemize}

\subsection{计算题答题规范}

\textbf{通用要求}：
\begin{itemize}
  \item 先写"解："，分步标号（1）（2）
  \item 写出原始公式（如 $F=ma$、$E=BLv$）再代入数值
  \item 代入数据时先写字母表达式，再代数字，最后写结果
  \item 注意单位，结果用科学记数法
  \item 最终答案建议单独一行或方框框出
  \item 矢量要说明方向
\end{itemize}

\textbf{各题型踩分关键}：

\[
\begin{array}{c|c}
\text{题型} & \text{踩分关键} \\ \hline
\text{力学综合} & \text{受力分析图 + 牛顿第二定律/动能定理/动量守恒}\\
\text{电学综合} & \text{电路分析 + 欧姆定律/电磁感应公式 + 能量转化}\\
\text{热学} & \text{理想气体状态方程 + 压强分析 + 热力学第一定律}\\
\text{光学} & \text{折射定律/全反射 + 几何关系}\\
\text{原子物理} & \text{质能方程/半衰期公式 + 核反应方程}
\end{array}
\]

\textbf{计算题步骤规范示例}（以力学综合为例）：

\begin{enumerate}
  \item 审题，画出受力分析图或运动过程图
  \item 明确研究对象，选择合适定理/定律
  \item 写出公式原型：$F_{\text{合}} = ma$
  \item 代入已知量：$mg\sin\theta - \mu mg\cos\theta = ma$
  \item 解方程得表达式：$a = g(\sin\theta - \mu\cos\theta)$
  \item 代数值计算，写结果（带单位）
  \item 检查合理性（方向、数值范围）
\end{enumerate}

\subsection{时间管理技巧}

\begin{itemize}
  \item \textbf{卡壳 3 分钟原则}：一道题思考超过 3 分钟没有进展，标记后跳过
  \item \textbf{首轮抢分}：先做所有能快速拿分的题
  \item \textbf{不空题}：选择题不会也猜一个（排除明显错误的选项后蒙）
  \item \textbf{计算题第一问必拿}：压轴题第一问通常是基础应用
  \item \textbf{最后 5 分钟}：检查答题卡填涂、实验题读数、计算题单位
\end{itemize}

\subsection{检查策略}

\begin{enumerate}
  \item \textbf{第 1 遍}（做完选择后）：检查单选题有无看错条件、矢量方向
  \item \textbf{第 2 遍}（做完实验后）：检查读数是否估读、电路连线是否正确
  \item \textbf{第 3 遍}（全部做完后）：检查计算题公式是否写对、单位是否遗漏
\end{enumerate}

\subsection{常见失分原因与对策}

\[
\begin{array}{c|c|c}
\text{失分原因} & \text{典型场景} & \text{对策} \\ \hline
\text{审题遗漏} & \text{漏看"光滑""不计电阻""缓慢"} & \text{圈出关键词再下笔} \\
\text{单位错误} & \text{cm 未换算成 m} & \text{统一使用 SI 单位制} \\
\text{方向混乱} & \text{正负号搞反} & \text{画图规定正方向} \\
\text{公式记混} & \text{左手/右手定则用反} & \text{左力右电} \\
\text{估算偏差} & \text{数量级算错} & \text{最后检查结果合理性} \\
\text{时间不足} & \text{在压轴题上纠结太久} & \text{3 分钟原则，果断跳过}
\end{array}
\]
